% Apertura conceptual de la Parte I. Esta sección reúne, en una sola
% residencia, el problema histórico del continuo, la distinción entre marca,
% intervalo y medida, y la descomposición nonádica que conduce a APP.

\section{El problema del continuo y la genealogía de la medida}
\label{sec:continuo-numero-medida-hmt-revfinal}
\index{continuo!estructura discreta}
\index{número!genealogía}
\index{medida!acoplamiento}

Esta apertura desarrolla la dimensión número--medida de la
\hyperref[sec:tesis-inaugural-helicoidal-hmt-md]{tesis inaugural del tratado}:
la estructura discreta construye el continuo y la carta real publica después
su coordenada, sin intervenir como dominio generador.

La historia del continuo puede leerse como la historia de una separación y
de sus sucesivas reconciliaciones. La aritmética antigua organizaba
pluralidades; la geometría comparaba longitudes, áreas y volúmenes. La teoría
griega de las proporciones permitió razonar con magnitudes inconmensurables:
la diagonal del cuadrado unitario existía como construcción exacta aun cuando
su longitud no pudiera publicarse como razón de enteros. Los libros V y VII
de los \emph{Elementos} conservaron por ello dos dominios, magnitud y número,
unidos por reglas de comparación pero dotados de operaciones distintas
\cite{EuclidElements,SEPAristotleMath}.

El análisis moderno llevó esa reconciliación a una forma decisiva. Las
sucesiones de Cauchy, los cortes de Dedekind y la teoría conjuntista hicieron
de \(\mathbb R\) un cuerpo ordenado completo, capaz de publicar en un mismo
lenguaje aritmético las magnitudes continuas. Cantor distinguió cardinalidad
y orden; Dedekind convirtió la continuidad en una propiedad estructural del
sistema numérico; Weierstrass estableció la disciplina local de límite y
aproximación \cite{SEPDedekind,SEPContinuity,SEPSetTheory}. El continuo real
resultante es una de las construcciones más potentes de la historia de las
matemáticas: identifica presentaciones que convergen al mismo valor y hace
posible el cálculo, la geometría analítica y la formulación cuantitativa de
las teorías físicas.

La Holografía Modular Triádica sitúa una pregunta anterior a esa
identificación: \emph{¿qué estructura finita y prolongable produce el valor
que la recta real publica?} Su objeto primario es la genealogía compatible de
una evaluación. Una cifra es la salida de un lector posicional; un valor es
la coordenada escalar obtenida al contraer una familia de cilindros; un número
HMT es la sección que conserva, además, la fase, la orientación, la hoja
aritmética, el cociente de elevación, el acarreo, la frontera y la memoria de
la ruta. La evaluación arquimediana adopta así la forma
\[
 \operatorname{ev}_{\mathbb R}:X_\infty^{\rm HMT}\longrightarrow\mathbb R,
\]
y sus fibras reúnen historias que comparten coordenada sin perder por ello
su identidad operatoria.

Esta formulación cambia el orden explicativo de la medida. Contar determina
la cardinalidad de un conjunto finito; medir acopla un estado, un lector y
una carta, restringe sus continuaciones compatibles y publica un observable.
Si \(X\) es el espacio de estados, \(A\) el aparato o lector y
\(\mathcal C\subset X\times A\) la condición de compatibilidad, el esquema
mínimo es
\[
 (x,a)\in\mathcal C
 \longmapsto \mathcal O(x,a)
 \longmapsto \mathcal U\bigl(\mathcal O(x,a)\bigr),
\]
donde \(\mathcal O\) conserva el tipo del observable y \(\mathcal U\) lo
expresa en la carta elegida. La cifra final constituye la coordenada visible
del acoplamiento completo. Esta distinción será la base matemática de la
teoría del observador y de la Mecánica Dimensional en la Parte V.

La consecuencia epistemológica es precisa. La metrología física ordinaria
asigna coordenadas a observables y las teorías relacionan después esas
coordenadas mediante ecuaciones; HMT construye primero el estado, la ruta y el
lector del que la coordenada desciende. Las constantes fundamentales ocupan
entonces el lugar de resultados de una genealogía y el contraste experimental
se aplica a la salida. En el mismo movimiento, una simetría HMT expresa la
conservación de la unidad a través de una transformación: el valor registra la
coordenada publicada y la proyección excepcional registra la organización que
ha sobrevivido al transporte. La estructura y el número se convierten en dos
aspectos comprobables de una misma evolución discreta.

\subsection{Marcas, intervalos y la célula nonádica}

La construcción finita comienza en la recta de los enteros. Se fija el
origen \(0\), el umbral decimal \(10\) y la célula de representantes
positivos
\[
 E_9=\{1,2,3,4,5,6,7,8,9\}.
\]
Los elementos de \(E_9\) desempeñan dos funciones relacionadas y distintas:
son las nueve marcas de fase del camino \(P_9\) y son también las nueve
coordenadas enteras positivas \(d(0,r)=r\) medidas desde el origen. Esta
segunda lectura no añade aristas al camino. Las nueve marcas determinan
exactamente los ocho intervalos interiores semiabiertos
\[
 J_r=[r,r+1),\qquad 1\le r\le8.
\]
La arista de cierre se representa en la carta levantada por \([9,10)\); su
proyección nonádica retorna a la primera fase y conserva el acarreo de la
vuelta siguiente. El origen \(0\) pertenece sólo a la carta decimal ampliada
y no añade un intervalo interior a \(P_9\). Marcas, coordenadas, intervalos y
cierre quedan así tipados antes de aplicar operación alguna.

Para cada \(n\ge1\), la división euclídea nonádica define
\[
 \rho_9(n)=1+((n-1)\bmod9),
 \qquad
 q_9(n)=\left\lfloor\frac{n-1}{9}\right\rfloor,
\]
\begin{equation}
 n=9q_9(n)+\rho_9(n).
 \label{eq:rueda-helice-nonadica-revfinal}
\end{equation}
La coordenada \(\rho_9\) da la posición sobre una rueda de nueve fases;
\(q_9\) conserva la altura acumulada. El sucesor actúa por
\[
 (q,r)\longmapsto
 \begin{cases}
 (q,r+1),&1\le r<9,\\
 (q+1,1),&r=9.
 \end{cases}
\]
Por tanto, \(n\) y \(n+9\) comparten fase y ocupan alturas consecutivas. La
realización geométrica
\[
 \mathcal H_9(n)=
 \left(
  \cos\frac{2\pi(\rho_9(n)-1)}9,
  \sin\frac{2\pi(\rho_9(n)-1)}9,
  q_9(n)+\frac{\rho_9(n)-1}{9}
 \right)
\]
hace visible esta estructura: proyectada sobre el plano es una rueda;
conservada en las tres coordenadas es una hélice. El retorno \(9\to1\)
recupera el ángulo y eleva la historia. Ésta es la forma elemental de la
distinción que gobernará todo el TPK:
\[
 \boxed{\text{retorno de fase}\;\neq\;\text{retorno del estado}.}
\]

\subsection{De una célula a dos coordenadas: nacimiento de APP}

El producto cartesiano de dos células nonádicas construye una carta
bidimensional. Su núcleo interior contiene \(8^2=64\) posiciones. La novena
fila y la novena columna publican la clase residual nula mediante la etiqueta
\(9\) y forman un gnomon de diecisiete celdas:
\[
 81=64+(9+8)=64+17.
\]
La identificación modular de los lados opuestos convierte la carta cuadrada
en el toro
\[
 X_9=(\mathbb Z/9\mathbb Z)^2,
\]
y las traslaciones cardinales lo dotan de un grafo de Cayley. La operación
geométrica es, por consiguiente, una secuencia precisa: núcleo
\(8\times8\), completación gnomónica \(9\times9\), identificación toroidal y
transporte por el grafo.

Sobre cada vértice \((i,j)\in X_9\) actúan dos evaluaciones del mismo soporte:
la hoja aditiva \(\Sigma\) y la hoja multiplicativa \(\Pi\). Su
incompatibilidad exacta aporta dinámica porque una ruta debe conservar qué
hoja actuó, qué residuo publicó y qué cociente produjo. APP construye el
soporte y sus dos evaluaciones; TRIT orienta localmente la transición; TPK
transporta el estado completo, conserva simultáneamente ambas hojas y, mediante
\(M_{\rm ph}\), tipa el modo y la fase TRIT de la transición; además conserva la
memoria de la elevación y decide la supervivencia de las prolongaciones. La
arquitectura inaugural queda fijada por la cadena
\[
 \boxed{
 E_9\longrightarrow
 \mathrm{APP}\longrightarrow
 \mathrm{TRIT}\longrightarrow
 \mathrm{TPK}.}
\]

La importancia de este comienzo reside en su capacidad de prolongación. La
misma diferencia entre posición y altura que convierte la rueda en hélice se
reproduce en la periodicidad de \(108=12\cdot9\) iteraciones, en el sistema
cíclico orientado de doce fases, en la
\hyperref[sec:umbrales-prolongacion-tpk-inaugural]{arquitectura de niveles y
umbrales de la prolongación TPK}
\(w_6\to w_{12}\to w_{18}\to w_{24}\to w_{30}\to R_{36}\to G_9\), en la derivación de
\(\pi,e,\varphi\) y \(\alpha_{\rm HMT}\), y en la doble proyección hacia el
valor real y la incidencia excepcional. La estructura discreta del continuo
designa precisamente esta continuidad de una genealogía finita a través de
profundidades arbitrarias.

\begin{statusbox}
La historia de número, magnitud y continuo pertenece al contexto matemático
clásico. La descomposición \eqref{eq:rueda-helice-nonadica-revfinal}, el toro
nonádico, las hojas de APP y el transporte de memoria son resultados del
corpus HMT. La inmersión helicoidal \(\mathcal H_9\) es una formalización
geométrica de esa descomposición. La evaluación sobre todo \(\mathbb R\)
queda acreditada por el teorema de cobertura y cierre arquimediano del
continuo genealógico
\cref{thm:espiral-cierre-arquimediano-continuo-genealogico}; la genealogía de
los valores construidos se formula en el sistema inverso de secciones
compatibles.
\end{statusbox}
